\documentclass{article}
\usepackage{amsmath, amssymb, amsthm}

% Section headings formatted as "1. Problem 4.12"
\makeatletter
\renewcommand{\@seccntformat}[1]{\csname the#1\endcsname.\enspace}
\makeatother

\newtheorem{lemma}{Lemma}
\newtheorem{theorem}{Theorem}
\theoremstyle{definition}
\newtheorem*{problem}{Problem}

\title{MATH 327 Homework 2 Fall 2026}
\author{Yazan Ismail}
\date{10/18/2026}

\begin{document}
\maketitle

%======================================================================
\section{Problem 4.12}

\begin{problem}
Let $\mathbb{I} = \mathbb{R} \setminus \mathbb{Q}$ denote the set of irrational numbers. Prove that if $a, b \in \mathbb{R}$ with $a < b$, then there exists $x \in \mathbb{I}$ such that $a < x < b$. (Hint: first show that $r + \sqrt{2}$ is irrational for every $r \in \mathbb{Q}$.)
\end{problem}

\begin{lemma}\label{lem:r-sqrt2}
Let $r \in \mathbb{Q}$. Then $r + \sqrt{2} \notin \mathbb{Q}$.
\end{lemma}

\begin{proof}
Suppose for the sake of contradiction that $r + \sqrt{2} \in \mathbb{Q}$. Thus, $r + \sqrt{2} = \frac{p}{q}$ for some $p \in \mathbb{Z}$, $q \in \mathbb{N}$, and $r = \frac{w}{v}$ for $w \in \mathbb{Z}$ and $v \in \mathbb{N}$. Hence, we have:
\begin{align*}
\frac{w}{v} + \sqrt{2} &= \frac{p}{q}, \\
\sqrt{2} &= \frac{p}{q} - \frac{w}{v}, \\
\sqrt{2} &= \frac{pv - wq}{qv},
\end{align*}
where $(pv - wq), (qv) \in \mathbb{Z}$ under the additive and multiplicative closure of $\mathbb{Z}$.

Hence, a contradiction, and thus $r + \sqrt{2} \notin \mathbb{Q}$.
\end{proof}

\begin{theorem}
If $a, b \in \mathbb{R}$ s.t.\ $a < b$, then there exists $x \in \mathbb{R} \setminus \mathbb{Q}$ s.t.\ $a < x < b$.
\end{theorem}

\begin{proof}
Let $a, b \in \mathbb{R}$ s.t.\ $a < b$. Thus, we have:
\begin{align*}
a &< b, \\
0 &< b - a, \tag{1} \\
0 &< \frac{b-a}{2}, \\
a &< \frac{b-a}{2} + a
\end{align*}
and we have
\begin{align*}
a &< b, \\
a + b &< 2b, \\
b - a &< 2b - 2a, \\
\frac{b-a}{2} &< b - a, \\
\frac{b-a}{2} + a &< b.
\end{align*}
Thus, we have
\begin{equation*}
a < \frac{b-a}{2} + a < b \tag{2}
\end{equation*}
via transitivity.

Since $\frac{b-a}{2} > 0$ by inequality (1), and since $\sqrt{2} > 0$, then by the Archimedean Property, $\exists\, n \in \mathbb{N}$ s.t.\ $n\left(\frac{b-a}{2}\right) > \sqrt{2}$,
\begin{equation*}
\frac{b-a}{2} > \frac{\sqrt{2}}{n}. \tag{3}
\end{equation*}
Let $n \in \mathbb{N}$ be such $n$.

Since $\sqrt{2}, n > 0$, then $\frac{\sqrt{2}}{n} > 0$ and hence:
\begin{equation*}
\frac{\sqrt{2}}{n} + a > a \tag{4}
\end{equation*}
Via transitivity of inequalities (2), (3), and (4) we have:
\begin{align*}
\frac{b-a}{2} + a &> \frac{\sqrt{2}}{n} + a > a, \\
b &> \frac{\sqrt{2}}{n} + a > a,
\end{align*}
and by Lemma~\ref{lem:r-sqrt2}, $\frac{\sqrt{2}}{n} + a \notin \mathbb{Q}$.

Hence, $\frac{\sqrt{2}}{n} + a \in \mathbb{I}$ s.t.\ $a < \frac{\sqrt{2}}{n} + a < b$.
\end{proof}

%======================================================================
\clearpage
\section{Problem 4.14(b)}

\begin{problem}
Let $A, B \subseteq \mathbb{R}$ s.t.\ $A$ and $B$ are nonempty and bounded. Let $A + B$ be the set of all sums $a + b$ where $a \in A$ and $b \in B$. Prove $\inf(A+B) = \inf(A) + \inf(B)$.
\end{problem}

\begin{proof}
Suppose that $\inf(A) + \inf(B) \neq \inf(A+B)$. Thus, $\inf(A) + \inf(B)$ is not a lower bound for $A+B$, or $\exists$ a lower bound for $A+B$ less than $\inf(A) + \inf(B)$.

\medskip
\underline{Case 1.} Suppose for the sake of contradiction that $\inf(A) + \inf(B)$ is not a lower bound for $A+B$. Thus, $\exists\, a + b \in A + B$ s.t.\ $a + b < \inf(A) + \inf(B)$. Since $a \geq \inf(A)$ $\forall a \in A$, then $0 \geq \inf(A) - a$. Hence,
\[
b < (\inf(A) - a) + \inf(B) \leq \inf(B),
\]
where $b < \inf(B)$ contradicts the def.\ of an infimum, and thus, $\inf(A) + \inf(B)$ is a lower bound for $A+B$.

\medskip
\underline{Case 2.} Suppose for the sake of contradiction that $\inf(A) + \inf(B)$ is not the greatest lower bound for $A+B$. Thus, $\forall\, a + b \in A + B$,
\[
a + b \geq \inf(A+B) > \inf(A) + \inf(B).
\]
Since $\inf(A+B) > \inf(A) + \inf(B)$, then there exists $x \in \mathbb{R}$ s.t.
\[
\inf(A) + \inf(B) < x < \inf(A+B) \leq a + b.
\]
Thus, we have:
\begin{align*}
\inf(A) + \inf(B) &< x < a + b, \\
\inf(A) + (\inf(B) - b) &< x - b < a,
\end{align*}
where $\inf(B) \leq b$, or $\inf(B) - b \leq 0$.

Suppose $\inf(B) = b$, or $\inf(B) - b = 0$. Thus, we have:
\[
\inf(A) + (\inf(B) - b) = \inf(A) < x - b < a,
\]
where $x - b$ is a greater lower bound for $A$ than $\inf(A)$, which is a contradiction.

Hence, we're left with $\inf(B) < b$. Thus, we have:
\[
\inf(B) < \frac{b - \inf(B)}{2} + \inf(B) < b,
\]
where $\frac{b - \inf(B)}{2} + \inf(B)$ is a greater lower bound for $B$ than $\inf(B)$, and we've thus exhausted all options with contradictions.

Hence, $\inf(A+B) = \inf(A) + \inf(B)$.
\end{proof}

%======================================================================
\clearpage
\section{Problem 8.2(c)}

\begin{problem}
Determine the limit of $c_n = \frac{4n+3}{7n-5}$ and prove that claim.
\end{problem}

\begin{theorem}
$\lim\left[\frac{4n+3}{7n-5}\right] = \frac{4}{7}$.
\end{theorem}

\begin{proof}
Let $c_n = \frac{4n+3}{7n-5}$. Let $\varepsilon > 0$. Let $N \in \mathbb{N}$ s.t.\ $\frac{1 + 35\varepsilon}{49\varepsilon} < N$.

Thus, we have:
\begin{align*}
\frac{1 + 35\varepsilon}{49\varepsilon} &< N, \\
\frac{1 + 35\varepsilon}{\varepsilon} &< 49N, \\
\frac{1}{\varepsilon} + 35 &< 49N, \\
\frac{1}{\varepsilon} &< 49N - 35, \\
\frac{1}{49N - 35} &< \varepsilon,
\end{align*}
and since $N \in \mathbb{N}$, then $49N - 35 \geq 0$. Hence, $\left|\frac{1}{49N - 35}\right| = \frac{1}{49N - 35}$.

Thus, we have:
\begin{align*}
\left|\frac{1}{49N - 35}\right| &< \varepsilon, \\
\left|\frac{(21 - 20) + (28N - 28N)}{49N - 35}\right| &< \varepsilon, \\
\left|\frac{(28N + 21) - 4(7N + 5)}{49N - 35}\right| &< \varepsilon, \\
\left|\frac{7(4N + 3) - 4(7N - 5)}{7(7N - 5)}\right| &< \varepsilon, \\
\left|\frac{4N + 3}{7N - 5} - \frac{4}{7}\right| &< \varepsilon.
\end{align*}
Thus, for an arbitrary $\varepsilon > 0$, there exists $N \in \mathbb{N}$ s.t.\ $n \geq N$ implies:
\[
\left|\frac{4N + 3}{7N - 5} - \frac{4}{7}\right| < \varepsilon.
\]
Hence, $\lim\left[\frac{4n+3}{7n-5}\right] = \frac{4}{7}$.
\end{proof}

%======================================================================
\clearpage
\section{Problem 8.4}

\begin{problem}
Let $(t_n)$ be a bounded sequence, i.e., there exists $M$ s.t.\ $|t_n| \leq M$ for all $n$, and let $(s_n)$ be a sequence s.t.\ $\lim s_n = 0$. Prove that $\lim(s_n t_n) = 0$.
\end{problem}

\begin{proof}
Let $s_n$ be a sequence s.t.\ $\lim[s_n] = 0$. Let $(t_n)$ be a bounded sequence s.t.\ $|t_n| \leq M$ for all $n \in \mathbb{N}$. Suppose for the sake of contradiction that $s_n t_n$ diverges.

Thus, there exists $\varepsilon > 0$ s.t.\ $\forall N \in \mathbb{N}$, $n \geq N$ and $|s_n t_n| \geq \varepsilon$.

Let $\varepsilon_2 > 0$ be such that
\begin{align*}
|s_n| M &\geq \varepsilon_2, \\
|s_n| &\geq \frac{\varepsilon_2}{M}.
\end{align*}
Let $\varepsilon_1 = \frac{\varepsilon_2}{M}$. Since $\varepsilon_2 > 0$ and $0 \leq |t_n| \leq M$, then $\frac{\varepsilon_2}{M} \geq 0$.

Thus, converging to $0$ implies that $\forall \varepsilon > 0$, $\exists N \in \mathbb{N}$ s.t.\ $n \geq N$ implies $|s_n| < \varepsilon$.

However, $\varepsilon_2$ is an $\varepsilon > 0$ s.t.\ $|s_n| \geq \varepsilon$. Hence, a contradiction.
\end{proof}

%======================================================================
\clearpage
\section{Problem 8.8(a)}

\begin{problem}
Prove $\lim\left[\sqrt{n^2 + 1} - n\right] = 0$.
\end{problem}

\begin{proof}
Let $\varepsilon > 0$. By the Archimedean Property, there exists $n \in \mathbb{N}$ s.t.
\[
\frac{1 - \varepsilon^2}{2\varepsilon} < n.
\]
Let $N$ be such $n$. Thus, we have:
\begin{align*}
\frac{1 - \varepsilon^2}{2\varepsilon} &< N, \\
\frac{1 - \varepsilon^2}{\varepsilon} &< 2N, \\
\frac{1}{\varepsilon} - \varepsilon &< 2N, \\
\frac{1}{\varepsilon} &< 2N + \varepsilon, \\
1 &< 2N\varepsilon + \varepsilon^2, \\
N^2 + 1 &< N^2 + 2N\varepsilon + \varepsilon^2, \\
\sqrt{N^2 + 1} &< \varepsilon + N,
\end{align*}
since $\varepsilon, N > 0$ so $\sqrt{N^2 + 1} > 0$ and hence $\sqrt{N^2 + 1} < -\varepsilon - N$ is a contradiction.

Thus, we have:
\begin{align*}
\sqrt{N^2 + 1} &< \varepsilon + N, \\
\sqrt{N^2 + 1} - N &< \varepsilon, \\
\sqrt{N^2\left(1 + \frac{1}{N^2}\right)} - N &< \varepsilon, \\
N\sqrt{1 + \frac{1}{N^2}} - N &< \varepsilon,
\end{align*}
where $\sqrt{1 + \frac{1}{N^2}} > 1$ $\forall N \in \mathbb{N}$. Hence, $N\sqrt{1 + \frac{1}{N^2}} > N$ and thus $\sqrt{N^2 + 1} - N > 0$.

Therefore,
\[
\sqrt{N^2 + 1} - N = \left|\sqrt{N^2 + 1} - N\right| < \varepsilon,
\]
where $\forall \varepsilon > 0$, we've shown there exists $N \in \mathbb{N}$ s.t.\ $n \geq N$ implies
\[
\left|\sqrt{N^2 + 1} - N\right| < \varepsilon, \quad \text{or} \quad \left|\left(\sqrt{N^2 + 1} - N\right) - 0\right| < \varepsilon.
\]

Hence, $\lim\left[\sqrt{n^2 + 1} - n\right] = 0$.
\end{proof}

%======================================================================
\clearpage
\section{Problem 6}

\begin{problem}
Prove that if $(s_n)$ converges, then so does $(|s_n|)$. Show by an example that the converse of this statement does not hold.
\end{problem}

\begin{proof}
Suppose $a_n$ converges to $a$. Let $\varepsilon_1$ be such an $\varepsilon$ s.t.\ $\frac{\varepsilon_1}{2} > 0$. Hence, there exists $N \in \mathbb{N}$ s.t.\ $n \geq N$ implies:
\begin{align*}
|a_n - a| &< \frac{\varepsilon_1}{2}, \\
2|a_n - a| &< \varepsilon_1, \\
|a_n - a| + |a_n - a| &< \varepsilon_1, \\
|a_n - a| &< \varepsilon_1 - |a_n - a|, \\
|a_n - a| &< \varepsilon_1 - \bigl|(a_n - |a_n|) + (|a_n| - a)\bigr| \\
&\leq \varepsilon_1 - \bigl|a_n - |a_n|\bigr| - \bigl||a_n| - a\bigr|, \tag{1}
\end{align*}
where inequality (1) holds via the triangle inequality. Thus, via transitivity, we have:
\begin{align*}
|a_n - a| &< \varepsilon_1 - \bigl|a_n - |a_n|\bigr| - \bigl||a_n| - a\bigr|, \\
|a_n - a| + \bigl|a_n - |a_n|\bigr| + \bigl||a_n| - a\bigr| &< \varepsilon_1.
\end{align*}
Since $|a_n - a|, \bigl|a_n - |a_n|\bigr| \geq 0$, then
\[
\bigl||a_n| - a\bigr| \leq |a_n - a| + \bigl|a_n - |a_n|\bigr| + \bigl||a_n| - a\bigr| < \varepsilon_1,
\]
and via transitivity, we have:
\[
\bigl||a_n| - a\bigr| < \varepsilon_1.
\]
Thus, for all $\varepsilon_1 > 0$, there exists $N \in \mathbb{N}$ s.t.\ $n \geq N$ implies $\bigl||a_n| - a\bigr| < \varepsilon_1$. Hence, $|a_n|$ converges and thus, $a_n$ converging implies $|a_n|$ converges.
\end{proof}

% TODO: Add an example showing the converse does not hold.

%======================================================================
\clearpage
\section{Problem 8.7(b)}

\begin{problem}
Prove $S_n = (-1)^n n$ diverges.
\end{problem}

\begin{proof}
Let $n \in \mathbb{N}$ and let $S_n = (-1)^n n$.

\medskip
\underline{Case 1.} Suppose $n = 2k$ for some $k \in \mathbb{Z}$. Thus,
\[
(-1)^{2k} \cdot n = \left[(-1)^k\right]^2 \cdot n,
\]
where $\forall x \in \mathbb{R}$, $x^2 \geq 0$. Thus since $n \geq 1$, then $\left[(-1)^k\right]^2 \cdot n \geq 0$, and hence, $(-1)^n = 1$ and $(-1)^n \cdot n = n$. Hence, $|(-1)^n \cdot n| = 1 \cdot n = n$.

\medskip
\underline{Case 2.} Suppose $n = 2j + 1$ for some $j \in \mathbb{Z}$. Thus,
\[
(-1)^n \cdot n = (-1)^{2j+1} \cdot n = (-1)^{2j} \cdot (-1) \cdot n,
\]
where $(-1)^{2j} \cdot n$ for $j \in \mathbb{Z}$ is equal to $n$. Thus,
\[
(-1)^n \cdot n = (-1)^{2j} \cdot (-1) \cdot n = (-1) \cdot n = -n.
\]
Hence, $|(-1)^n \cdot n| = |-n| = n$ since $n \geq 1$.

\medskip
Thus, $|S_n| = |(-1)^n n| = n$.

Since $|S_n| = n$ is a trivially divergent sequence, then by Problem 6 of this homework (HW 2), the contrapositive says:
\begin{center}
``If $|a_n|$ diverges, then $a_n$ diverges.''
\end{center}
Hence, since $|S_n| = n$ diverges, then $S_n = (-1)^n n$ diverges.
\end{proof}

%======================================================================
\clearpage
\section{Problem 8.9(a)}

\begin{problem}
Let $(s_n)$ be a convergent sequence. Show that if $s_n \geq a$ for all but finitely many $n$, then $\lim s_n \geq a$. (Per the assignment, you may assume the inequality holds for all $n$.)
\end{problem}

% TODO: Add your proof of 8.9(a) here.
% \begin{proof}
%
% \end{proof}

\end{document}
